\chapter{心智模型与形成性评价}\label{s:models}

教学的第一项任务，是弄清你的学习者是谁。
我们的方法是基于 Patricia Benner 等研究者的工作，\index{Benner, Patricia}
他们研究了护士如何从新手成长为专家~\cite{Benn2000}。
Benner 识别出认知发展的五个阶段，
大多数人都以相当一致的方式经历这些阶段。
就我们的目的而言，可以把这一进程简化为三个阶段：

\begin{description}

\item[\grefdex{g:novice}{新手}{新手}]
  不知道自己不知道什么，
  也就是还没有对问题领域建立起可用的心智模型。

\item[\grefdex{g:competent-practitioner}{熟手}{熟手}]
  拥有足以应付日常用途的心智模型。
  他们能在正常情况下以正常的努力完成常规任务，
  并且对自己知识的边界有一定认识
  （也就是知道自己不知道什么）。

\item[\grefdex{g:expert}{专家}{专家}]
  的心智模型包含例外和特例，
  这使他们能处理不寻常的情况。
  我们会在\chapref{s:memory}里更详细地讨论专长。

\end{description}

那么，\gref{g:mental-model}{心智模型}\emph{究竟}是什么？
顾名思义，
它是对某个问题领域中最重要的部分所做的简化表示，
其精确程度足以支撑解决问题。
一个例子是高中化学里用来表示分子的球—弹簧模型。
原子其实不是球，
它们的化学键其实也不是弹簧，
但这个模型能让人对化合物及其反应进行推理。
更精细的原子模型是中心一个小球（原子核），周围环绕着电子。
它也是错的，
但额外的复杂性让人能解释更多、解决更多问题。
（和软件一样，
心智模型永远没有完成：
它们只是被拿来用而已。）

给新手一下子抛出一堆事实是适得其反的，
因为他们还没有一个模型来安放这些事实。
事实上，
过早抛出太多事实，反而可能强化
他们东拼西凑起来的那套错误心智模型。
正如~\cite{Mull2007a}在一项针对理科生的视频教学研究中观察到的：

\begin{quote}

  学生在观看视频之前，对{\ldots}现象已有自己的想法。
  如果视频清晰、配图充分地呈现{\ldots}概念，
  学生会以为自己在学，但他们与媒体的互动深度还不足以
  意识到所呈现的内容和他们原有的知识并不一样{\ldots}
  不过还是有希望的。
  在视频里把学生常见的误解和{\ldots}概念一起呈现，
  已被证明能增加学习，
  因为它增加了学生观看时投入的心理努力。

\end{quote}

所以你在教新手时的目标，
应该是帮助他们建构一个心智模型，
好让他们有地方安放事实。
例如，
Software Carpentry 的\index{The Carpentries}
\hreffoot{http://swcarpentry.github.io/shell-novice/}{Unix shell 课程}
用三小时介绍了十五条命令。
也就是每十二分钟一条命令，
这看起来慢得像冰川移动，直到你意识到
这门课的真正目的并不是教这十五条命令：
而是教路径、
历史、
Tab 补全、
通配符、
管道、
命令行参数
和重定向。
只有新手理解了这些概念，具体的命令才有意义；
一旦他们理解了，
就能开始读手册页、
在网上搜索合适的关键词，
并判断搜到的结果有没有用。

新手和熟手之间的认知差异，
支撑着两类教学材料的差异。
\gref{g:tutorial}{教程}帮助刚入门的人建立心智模型；
而\gref{g:manual}{手册}
则帮助熟手填补知识上的空缺。
教程会让熟手沮丧，因为它进展太慢，
讲的都是显而易见的东西
（尽管对新手来说它们\emph{一点也不}显而易见）。
同样，
手册会让新手沮丧，因为它满是行话，而且\emph{不}解释。
这种现象叫做\gref{g:expertise-reversal}{专长逆转效应}~\cite{Kaly2003}，
也是你必须尽早决定
课程是为谁而讲的其中一个理由。

\begin{aside}{少数的例外}
  Unix 和 C 之所以流行，其中一个原因是
  \cite{Kern1978,Kern1983,Kern1988}
  竟然同时既是好教程\emph{又是}好手册。
  \cite{Fehi2008} 和~\cite{Ray2014} 是计算机领域仅有的几本做到这一点的书；
  即便反复重读，我也仍然不明白它们是怎么做到的。
\end{aside}

\seclbl{人们学到了吗？}{s:models-formative-assessment}

马克·吐温曾写道，\index{Twain, Mark}
「害你惹上麻烦的，不是你不知道的事，
而是你确信不疑、其实并非如此的事。」
因此，建构心智模型的功课之一，
就是把\emph{不}属于其中的东西清理掉。
大致说来，
新手的误解可以分为三类：

\begin{description}

\item[事实性错误]
  比如以为温哥华是不列颠哥伦比亚省的首府
  （其实是维多利亚）。
  这类错误通常很容易纠正。
  \index{误解!事实性错误}

\item[模型损坏]
  比如以为运动和加速度的方向必须一致。
  我们可以让新手推演一些例子，
  在这些例子里他们的模型会给出错误答案，从而纠正它。
  \index{误解!损坏的心智模型}

\item[根本信念]
  比如「世界只有几千年历史」，
  或者「某些人天生就比别人更擅长编程」~\cite{Guzd2015b,Pati2016}。
  这类错误往往和学习者的社会身份深深纠缠在一起，
  因此会抗拒证据、并给矛盾找理由。
  \index{误解!根本信念}

\end{description}

当教师能在授课过程中识别并清除学习者的误解时，人们学得最快。
这叫做\gref{g:formative-assessment}{形成性评价}，
因为它在教学进行的同时形成（或塑造）教学。
学习者不会在形成性评价中及格或不及格；
相反，
它让教师和学习者都得到反馈，知道目前做得怎么样、
接下来该把重点放在哪里。
例如，
音乐老师可能让学习者非常缓慢地弹一段音阶，来检查他的呼吸。
学习者由此知道自己呼吸对不对，
而老师则得到反馈，知道他刚才的讲解有没有讲清楚。

\begin{aside}{总结一下}
  与形成性评价相对的是
  \gref{g:summative-assessment}{总结性评价}，
  它发生在一门课程结束时。
  总结性评价就像驾照考试：
  它告诉学习者是否掌握了这个主题，
  也告诉老师他的课是否成功。
  理解两者区别的一种方式是：
  厨师做菜时尝味道是形成性评价，
  而客人吃到端上来的菜则是总结性评价。

  遗憾的是，
  学校已经让大多数人相信所有评价都是总结性的，
  也就是只要感觉像是考试，
  考得差就会对你不利。
  让形成性评价显得不那么正式，有助于减轻这种焦虑；
  依我的经验，
  用在线测验、答题器或别的东西，反而会加重焦虑，
  因为今天大多数人相信自己在网上做的任何事都会被监视和记录。
\end{aside}
  
为了在教学中派上用场，
形成性评价必须能快速实施
（以免打断课的节奏），
而且要有唯一明确的正确答案
（这样才能用于群体）。\index{形成性评价!要求}
用得最广的一类形成性评价大概是
多项选择题（MCQ）。\index{多项选择题}
很多老师看不上它，
但如果设计得当，
它能揭示的远不止某人是否知道某些具体事实。
例如，
假设你在教孩子们做多位数加法~\cite{Ojos2015}，
你给他们这样一道 MCQ：

\begin{quote}
  37 + 15 等于多少？\\
  a) 52\\
  b) 42\\
  c) 412\\
  d) 43
\end{quote}

\noindent
正确答案是 52，
但其他答案提供了有价值的线索：

\begin{itemize}

\item
  如果孩子选了 42，
  说明他根本不理解「进位」是什么意思。
  （他很可能是把 7+5 的答案写成 12，
  然后用 3+1 得到的 4 覆盖掉那个 1。）

\item
  如果他选了 412，
  说明他把每一列数字当成了独立的问题。
  这仍然不对，
  但错的原因不一样。

\item
  如果他选了 43，说明他知道要把 1 进位，
  却是把它进回到了原来那一列。
  这又是另一种错误，
  需要老师给出另一种有针对性的解释。

\end{itemize}

这些错误答案每一个都是具有\gref{g:diagnostic-power}{诊断力}的
\gref{g:plausible-distractor}{看似合理的干扰项}。
干扰项是错误答案或不那么理想的答案；
「似真」的意思是它看起来像是正确答案，
而「诊断力」的意思是每个干扰项都能帮老师弄清
接下来该给这位学习者讲什么。

形成性评价中答案的分布，会引导你下一步怎么做。
如果班里足够多的人答对了，你就往下走。
如果班里大多数人都选了同一个错误答案，
你就应该回去，纠正那个干扰项所指向的误解。
如果他们的答案均匀地分散在几个选项上，那很可能只是在猜，
所以你应该退回去，换一种方式重新解释这个想法。
（原封不动地重复同样的解释大概没用，
这也是那么多视频课程在教学上效果不佳的原因之一。）

要是班里大多数人选了正确答案，
但少数人选错了呢？
这时，
你就得决定：是把时间花在让少数人跟上，
还是让大多数人保持投入更重要。
不管你多么努力、用什么教学方法，
你都不可能总让每个人都得到他们需要的东西；
作为老师，做出这个判断是你的责任。

\begin{aside}{错误答案是打哪儿来的？}
  要想出似真的干扰项，
  可以回想上一次教这个主题时，
  学习者问过哪些问题、遇到过哪些困难。
  如果你以前没教过，
  就回想一下你自己有过的误解，
  问问同事的经验，
  或者看看你这个领域的历史：
  如果五十年前大家都以某种方式误解了你的学科，
  那么今天你的许多学习者很可能仍然那样误解它。
  你也可以在课上问一些开放性问题，
  收集关于以后要讲的内容有哪些误解，
  或者去看看 \hreffoot{http://www.quora.com}{Quora}
  或 \hreffoot{https://stackoverflow.com/}{Stack Overflow} 之类的问答网站，
  看看在别处学这门学科的人都被什么搞糊涂了。
\end{aside}

制作形成性评价能让你的课更好，
因为它迫使你思考学习者的心智模型。
依我的经验，
一旦这样做，我就会自动把课写成能覆盖最可能出现的空缺和错误。
所以，即使形成性评价没被用上，它也是值得的
（当然，用上它教学会更有效）。

MCQ 并不是唯一一类形成性评价：
\chapref{s:exercises}描述了其他几种既快又无歧义的练习。
不管你选哪种，
都应该每 10—15 分钟做一两分钟的事，
以确认你的学习者真的在学。
这个节奏并不是基于某种内在的注意力极限：~\cite{Wils2007}
几乎不支持那个被反复念叨的说法，
即学习者只能集中注意力 10—15 分钟。
相反，
这个指导原则是为了保证：如果相当多的人掉队了，
你只需要把这门课的很小一部分再讲一遍。
频繁的形成性评价还能让学习者保持投入，
尤其是当它涉及小组讨论时
（\secref{s:classroom-peer}）。

形成性评价也可以用在一门课\emph{之前}。
如果你用一道 MCQ 开始上课，而所有人都答对了，
你就可以不必去讲学习者已经知道的东西。
这种\gref{g:active-teaching}{主动教学}
让你有更多时间专注于他们不知道的东西。
它也向学习者表明，你尊重他们的时间、不至于浪费它，
这有助于激发动机（\chapref{s:motivation}）。

\begin{aside}{概念清单}
  只要有足够的数据，
  MCQ 可以被做得精确得惊人。
  最著名的例子是
  \grefdex{g:concept-inventory}{力概念清单}{概念清单}~\cite{Hest1992}，
  它评估对基础牛顿力学的理解。
  通过访谈大量受访者、
  把他们的误解与对错答案的模式关联起来，
  再改进题目，
  它的作者们构建出一个能精确定位具体误解的诊断工具。
  研究者随后就能用这个工具来衡量教学方法变化的效果~\cite{Hake1998}。

  Tew 等人开发并验证了一套
  与编程语言无关的入门编程评估~\cite{Tew2011}；
  \cite{Park2016} 做了重复验证，
  而~\cite{Hamo2017} 正在开发一套针对递归的概念清单。
  然而，
  建立这样的工具成本极高，
  而且学习者上网搜答案的能力，正越来越威胁到它们的有效性。
\end{aside}

把形成性评价融入课堂，只需要一点点准备和练习。
给学习者发彩色或带编号的卡片，让他们能同时回答一道 MCQ
（而不是一个个举手），
把其中一个选项设成「我完全不知道」，
并鼓励他们在作答前和邻座讨论几秒钟，
这些都有助于确保你的教学节奏不被打断。
\secref{s:classroom-peer}描述了一种强大、
基于证据的教学方法，就建立在这些简单的想法之上。

\begin{aside}{幽默}
  老师有时会在 MCQ 里放一些本意是搞笑的答案，比如「我的鼻子！」，
  尤其是给年龄较小的学习者出的题。
  然而，
  这些答案对了解学习者的误解毫无帮助，
  而且大多数人其实并不觉得它们好笑。
  一般而言，
  只有当你在第三次重读时仍然觉得某个笑话好笑，
  你才该把它放进课程里。
\end{aside}

一门课的形成性评价应当为它的总结性评价做好准备：
任何人都不应该遇到一道考试题目，而教学却没有为他做好准备。
这并不意味着你绝不能在考试里放新的题型，
但如果你这么做，
就应该事先让学习者练习过处理新颖的问题。
\chapref{s:process}会更深入地探讨这一点。

\seclbl{观念机器}{s:models-notional}

\gref{g:computational-thinking}{计算思维}这个词被到处滥用，
部分原因是人们可以一致同意它很重要，心里想的却是非常不同的东西。
与其争论它包含或不包含什么，
不如想一想你希望学习者理解的\gref{g:notional-machine}{观念机器}~\cite{DuBo1986}。
根据~\cite{Sorv2013}，
观念机器：

\begin{itemize}

\item
  是对计算机硬件
  以及程序运行时环境其他方面的一种理想化抽象；

\item
  能让人描述程序的语义；
  并且

\item
  能正确反映程序执行时的行为。

\end{itemize}

\noindent
例如，我为 Python 建立的观念机器是这样的：\index{Python}

\begin{enumerate}

\item
  运行中的程序住在内存里，
  内存分为调用栈和堆两部分。

\item
  数据的内存总是从堆里分配的。

\item
  每份数据都存放在一个由两部分组成的结构里。
  第一部分说明数据是什么类型，
  第二部分才是真正的值。

\item
  布尔值、数字和字符串一旦创建就不再被修改。

\item
  列表、集合和其他容器存放的是对其他数据的引用，
  而不是直接存放那些值。
  它们在创建之后可以被修改，
  也就是列表可以被扩展，集合可以加入新值。

\item
  当代码被载入内存时，
  Python 把它转换成一串指令，
  这些指令和其他数据一样被存储。
  这就是为什么可以把函数赋给变量、
  并作为参数传递。

\item
  代码执行时，
  Python 逐条走过这些指令，
  依次做每一条告诉它做的事。

\item
  有些指令让 Python 读取数据、
  做计算、
  并创建新数据。
  另一些指令控制 Python 执行哪些指令，
  循环和条件就是这样工作的。
  还有一类指令告诉 Python 去调用一个函数。

\item
  调用函数时，
  Python 把一个新栈帧压入调用栈。

\item
  每个栈帧存放变量的名字和对数据的引用。
  函数参数不过是另一种变量。

\item
  用到某个变量时，
  Python 先在栈顶的帧里找它。
  如果不在那里，就去最底下（全局）的帧里找。

\item
  函数结束时，
  Python 抹掉它的栈帧，跳回到
  调用函数之前正在执行的指令。
  如果没有「之前」，
  程序就结束了。

\end{enumerate}

我教 Python 时用的就是这种漫画版的现实。
经过大约 25 小时的授课和 100 小时的自主练习，
我预期大多数学习者会建立起一个心智模型，
包含这里的大部分或全部特征。

\seclbl{练习}{s:models-exercises}

\exercise{你的心智模型}{思考—结对—分享}{15}

你用来理解自己工作的一种心智模型是什么？
写几句话描述它，并对同伴的模型给出反馈。
做完之后，
请几个人把他们的模型分享给全班。
大家对一个心智模型是什么有共识吗？
有可能给出一个精确的定义吗，
还是说这个概念之所以有用，恰恰因为它本身是模糊的？

\exercise{新手的一些症状}{全班}{5}

说新手对某个领域没有心智模型，
并不等于说他们完全没有心智模型。
新手往往靠类比和猜测来推理，
从其他看起来表面上相似的领域借来一鳞半爪的心智模型。

这样做的人常说一些
\hreffoot{https://en.wikipedia.org/wiki/Not\_even\_wrong}{连错都算不上}的话。
作为全班一起讨论，新手的其他一些症状是什么。
一个人做了什么、说了什么，会让你判定他在某个领域是新手？

\exercise{给新手的心智模型建模}{两人一组}{20}

针对你教过或打算教的一个主题，设计一道多项选择题，
并解释它每个干扰项的诊断力
（也就是每个干扰项分别想识别出哪种误解）。

做完之后，和同伴交换 MCQ。
他的题目有歧义吗？
那些误解似真吗？
干扰项真的能测出这些误解吗？
有没有一些很可能存在的误解\emph{没有}被测到？

\exercise{把事情想透}{全班}{15}

好的形成性评价要求人去想透一个问题。
例如，
设想你把一块冰放进浴缸，然后把浴缸放满水、直到水漫到边缘。
冰融化后，
水位会上升（于是浴缸溢水）、
下降，
还是保持不变（\figref{f:models-bathtub}）？

\figpdf{figures/bathtub.pdf}{浴缸里的冰}{f:models-bathtub}{}

正确答案是水位保持不变：
冰排开的是与自身重量相等的水，
所以融化后正好填满它此前「挖」出的那个洞。
弄清其中缘由，有助于人们建立一个关于重量、体积和密度之间关系的心智模型~\cite{Epst2002}。

再描述一个你见过或用过的形成性评价，
它要求人把某件事想透，
从而发现他们推理中的缺陷。
做完之后，
向同伴解释你的例子，
并对他的例子给出反馈。

\exercise{另一种进阶}{个人}{15}

这种新手—熟手—专家的技能发展模型
有时叫做\hreffoot{https://en.wikipedia.org/wiki/Dreyfus\_model\_of\_skill\_acquisition}{德雷福斯模型}。
另一种常用的进阶是
\hreffoot{https://en.wikipedia.org/wiki/Four\_stages\_of\_competence}{能力四阶段}：

\begin{description}

\item[无意识的无能：]
  此人不知道自己不知道什么。

\item[有意识的无能：]
  此人意识到自己不知道某件事。

\item[有意识的有能：]
  此人学会了做某件事，
  但只有集中注意力时才能做到，
  可能仍需要把它拆成步骤。

\item[无意识的有能：]
  这项技能已经变成第二天性，
  此人可以反射性地完成它。

\end{description}

找出一个你分别处于这四个阶段的科目。
在你最常教的科目里，你的大多数学习者处于哪个阶段？
你想把他们带到哪个阶段？
这四个阶段与新手—熟手—专家的分类有什么关系？

\exercise{哪一种计算？}{个人}{10}

\cite{Tedr2008} 总结了计算领域的三种传统：

\begin{description}

\item[数学传统：]
  程序是算法的化身。
  它们要么正确要么不正确，
  也分效率高低。

\item[科学传统：]
  程序是对信息过程的或精确或不精确的模型，
  可以用科学方法来研究。

\item[工程传统：]
  程序像水坝和飞机一样是人造物，
  分有效可靠与否。

\end{description}

哪一种最符合你对计算的心智模型？
如果都不是，那你拥有的是什么模型？

\exercise{解释为什么不是这样}{两人一组}{5}

你的一个学习者认为，
他一个字符一个字符敲进去的文本，
和他们复制粘贴的完全相同的文本之间，
存在某种差别。
想想他为什么可能相信这一点，
或者可能发生过什么事让他产生这种印象，
然后假装你就是那个学习者，而你的同伴向你解释为什么并非如此。
交换角色再试一次。

\exercise{你此刻的模型}{全班}{5}

作为全班，
列出你们关于学习的心智模型的关键要素。
最重要的六七个概念是什么，它们之间又是什么关系？

\exercise{你的观念机器}{小组}{20}

以小组为单位，
写出你希望学习者用来理解程序如何运行的观念机器的描述。
像 Scratch 这样的积木式语言，其观念机器与 Python 的有什么不同？\index{Scratch}\index{Python}
电子表格的观念机器呢？
或者一个渲染网页时解释 HTML 和 CSS 的浏览器呢？

\exercise{只乐不学}{个人}{5}

多项研究表明，
教学评价与学习成效并不相关~\cite{Star2014,Uttl2017}，
也就是学习者对一门课的评价有多高，并不能预测他们记住了多少。
你有没有上过一门自己很享受、但其实什么也没学到的课？
如果有，是什么让它变得有趣？

\newpage
\section*{回顾}

\figpdfhere{figures/conceptmap-mental-models.pdf}{概念：心智模型}{f:models-concept-map}{}

\figpdfhere{figures/conceptmap-assessment.pdf}{概念：评价}{f:models-formative-assessment}{}
